% Documentation text for `screenwriter-doc.tex', version 1.0.1
% Paper size: A4
% J. A. Corbal, 2026

\documentclass[letter]{article}

\usepackage[T1]{fontenc}
%\usepackage[utf8]{inputenc}
\usepackage[english]{babel}
\usepackage{lmodern}
\usepackage{microtype}
\usepackage[colorlinks=true,
            linkcolor=blue,
            urlcolor=blue,
            citecolor=blue]{hyperref}

\hypersetup{
    pdftitle = {The screenwriter class},
    pdfauthor = {J. A. Corbal},
    pdfsubject = {Documentation for the screenwriter LaTeX class},
    pdfkeywords = {LaTeX, screenplay, screenwriter, class},
    %pdfcreator = {LaTeX},
}

\newcommand{\meta}[1]{<\textit{#1}>}

\title{The \textsf{screenwriter} class\thanks{This document was last
    revised and validated on \today.}}

\author{J.~A.~Corbal\thanks{\texttt{jacorbal [at] gmail [dot] com}}}

\date{Version~1.0.1 --- 2026-05-09}


\begin{document}%{{{

\maketitle

\begin{abstract}

    The \textsf{screenwriter} class provides a LaTeX document class for
    writing film and television screenplays.  It is based on John Pate's
    \texttt{screenplay} class, updated for current LaTeX distributions
    and extended with modular language support and better integration
    with \texttt{babel}.

\end{abstract}

\section{Introduction}%{{{

The \textsf{screenwriter} class is a \LaTeX{} document class designed to
typeset film and television screenplays.  It provides the conventional
screenplay layout (12\nobreak\thinspace pt monospaced type, specific
margins, dialogue blocks, and sluglines) and a set of high-level
commands tailored to the way screenwriters actually work.

The class is based on John Pate's original \texttt{screenplay.cls}, but
has been updated and extended for current \LaTeX{} systems.  In
particular, it modernises package loading, separates language-specific
text into dedicated language definition files, and offers better
cooperation with the \texttt{babel} package.

Unlike generic article or book classes, \textsf{screenwriter} focuses on
the visual and structural conventions of screenplays: scene headings,
dialogue layout, transitions such as ``FADE IN:'' and ``FADE OUT:'', and
a standardised cover page.  The goal is to let authors concentrate on
the script itself, while \LaTeX{} takes care of consistent formatting.

The class is intended both for individual writers who like working with
\LaTeX, and for production environments where reproducible layouts and
multi-language support are important.

%}}}

\section{Installation}%{{{

This section describes a typical local installation in a personal \TeX{}
tree.  System-wide installation follows the same pattern, but requires
administrator privileges.

The distribution of \textsf{screenwriter} consists of the following
types of files:

\begin{itemize}\itemsep0em
  \item \texttt{screenwriter.cls}: the main document class.
  \item \texttt{screenwriter-lang-\meta{language}.ldf}: language
        definition files containing all language-dependent strings used
        by the class (e.g., ``FADE IN'', ``THE END'', ``INT.'', ``EXT.'').
  \item \texttt{screenwriter-doc.tex} \& \texttt{screenwriter-doc.pdf}:
        this documentation and its source.
  \item \texttt{LICENSE}: license file (LPPL\,1.3c).
  \item \texttt{README}: this file.
\end{itemize}

A typical \TeX{} directory structure might look like:

\begin{verbatim}
    screenwriter/
      LICENSE
      README
      doc/
        screenwriter-doc.tex
        screenwriter-doc.pdf
      tex/
        latex/
          screenwriter/
            screenwriter.cls
            lang/
              screenwriter-lang-english.ldf
              screenwriter-lang-spanish.ldf
              ...
\end{verbatim}

For a local installation, copy the \texttt{tex/latex/screenwriter/}
contents into your personal \TeX{} tree, for example:

\begin{center}
  \texttt{\$TEXMFHOME/tex/latex/screenwriter/}
\end{center}
where \texttt{\$TEXMFHOME} is usually \texttt{\~{}/texmf} on a standard
\TeX{}~Live installation.  The language definition files can be placed
in a \texttt{lang/} subdirectory as shown above; the class looks for
them as \texttt{lang/screenwriter-lang-\meta{language}.ldf}.

After copying the files, update your file name database if required by
your distribution (for example by running \texttt{mktexlsr} on \TeX{}
Live).  Once this is done, \LaTeX{} will be able to find the class and
language files automatically, and you can use in any document:

\begin{verbatim}
    \documentclass{screenwriter}
\end{verbatim}

%}}}

\section{Basic usage}%{{{

This section shows a minimal example and explains the main elements that
\textsf{screenwriter} provides out of the box.

\subsection*{A minimal document}

A very small English screenplay using the default layout might look as
follows:

\begin{verbatim}
    \documentclass{screenwriter}

    \title{My First Screenplay}
    \author{Screenwriter Name}

    \begin{document}
        \coverpage

        \fadein
        \intslug[DAY]{OFFICE}

        \begin{dialogue}{WRITER}
            This is a line of dialogue.
        \end{dialogue}

        \theend
    \end{document}
\end{verbatim}

The class sets up 12\nobreak\thinspace pt monospaced type, appropriate
margins and dialogue widths, and a standard page style.  The example
uses only a few of the commands provided, but already shows the typical
structure: cover page, opening transition, first scene heading,
dialogue, and a final ``THE END'' marker.

\subsection*{Basic structure}

In a typical screenplay document:

\begin{itemize}
  \item The preamble selects the class and language, using \\
        \verb|\documentclass[⟨options⟩]{screenwriter}|, and provides
        meta-data such as \verb|\title| and \verb|\author|.
  \item The document body starts with \verb|\coverpage| or \verb|\nicholl|
        to produce a title page.
  \item Each scene begins with a slugline, such as \verb|\intslug| or
        \verb|\extslug|.
  \item Dialogue is set within the \verb|dialogue| environment, with
        optional parentheticals and continuation marks handled by
        \verb|\dialbreak|.
  \item Transitions such as \verb|\fadein|, \verb|\fadeout|,
        \verb|\intercut| and \verb|\flashback| are used where needed.
  \item The script can end with \verb|\theend|, which inserts the
        language-dependent ``THE END'' marker centred on the page.
\end{itemize}

%}}}

\section{Options}%{{{

The \textsf{screenwriter} class accepts a small set of options that
control page layout, screenplay style, and language.

\subsection{Paper and layout options}

The class internally loads the standard \texttt{article} class with
12\nobreak\thinspace pt, one-column, one-sided settings, and then
adjusts margins and text width to match common screenplay practice.

Two paper size options are available:
\begin{description}
  \item[\texttt{letterpaper}] Use US Letter paper (default).  This is
      the usual choice for North American markets.
  \item[\texttt{a4paper}] Use A4 paper and adjust the text and dialogue
      widths accordingly.  This is often more convenient for European
      printers and readers.
\end{description}

The screenplay style can be selected with:
\begin{description}
  \item[\texttt{literary}] Literary (spec) screenplay style.  Scene
        headings (sluglines) are printed without scene numbers.  This is
        the default.
  \item[\texttt{technical}] Technical (numbered) style.  Each scene
        heading is automatically numbered, and the scene number is
        printed in the margins.  This is useful in production and
        shooting scripts.
\end{description}

Examples:
\begin{verbatim}
    \documentclass[a4paper]{screenwriter}
    \documentclass[letterpaper,technical]{screenwriter}
\end{verbatim}

\subsection{Language options}

The class separates language-dependent strings (such as ``FADE IN'',
``THE END'', ``INT.'', ``EXT.'', ``TITLE OVER'', etc.) into external
language definition files \verb|screenwriter-lang-⟨language⟩.ldf|.  The
language used by the class can be selected in two ways.

\subsubsection*{Explicit language selection}

You can explicitly choose the language via class options.  The following
languages are currently provided\footnote{The language files should
undergo a formal review to ensure accuracy and consistency.}:
\begin{center}

    \texttt{english}, \texttt{spanish}, \texttt{french},
    \texttt{german}, \texttt{italian}, \\
    \texttt{portuguese}, \texttt{galician}, \texttt{catalan},
    \texttt{esperanto}

\end{center}

For example:
\begin{verbatim}
    \documentclass[spanish]{screenwriter}
\end{verbatim}
tells the class to use Spanish strings and to therefore load the
language file \verb|lang/screenwriter-lang-spanish.ldf|, regardless of
whether the document uses \texttt{babel} or not.

\subsubsection*{Cooperation with \texttt{babel}}

If no explicit language option is given to \textsf{screenwriter}, but
the user loads \texttt{babel} with a language, the class will attempt to
detect the current \texttt{babel} language via \verb|\languagename| at
\verb|\begin{document}|.  It then tries to load the corresponding
language file:
\begin{verbatim}
    \documentclass{screenwriter}
    \usepackage[french]{babel}
\end{verbatim}

In this case, the class will look for
\verb|lang/screenwriter-lang-french.ldf| and use French strings for all
its internal labels.  If no matching language file is found, the class
falls back to English.

\paragraph{Note.}  The class \textsf{screenwriter} does \emph{not} load
\texttt{babel} on its own; users are free to choose \texttt{babel},
\texttt{polyglossia}, or no multilingual package at all.  The class only
mirrors the active language when possible.

%}}}

\section{Commands}%{{{

This section summarises the main commands and environments provided by
the \textsf{screenwriter} class.  They are grouped by function: title
and address information, sluglines (scene headings), dialogue, and
transitions.

\subsection{Title and cover page}

The class reuses the standard \LaTeX{} meta-data commands
\verb|\title| and \verb|\author|, and adds a few macros specific to
screenplays.

\begin{description}
  \item[] \verb|\title{<title>}|\quad  Sets the script title.  This is
    used on the cover page and may also appear on the first page of the
    script.
  \item[] \verb|\author{<name>}|\quad  Sets the author name shown
    beneath the title.
  \item[]\verb|\realauthor{<name>}|\quad  Sets the real author name,
    as printed in the address block on the cover when \verb|\author|
    refers to a pseudonym.  By default, this is the same as
    \verb|\author|.
  \item[]\verb|\address{<text>}|\quad  Sets the contact address that
    appears in the cover page address block.  The default is a generic
    ``Contact via Agency'' (or its language-dependent equivalent).
  \item[]\verb|\agent{<text>}|\quad  Sets the agent information for the
    cover page.  By default, this is left empty.
\end{description}

Once these values are set, you can produce a cover page using:
\begin{description}
  \item[] \verb|\coverpage|\quad  Typesets a standard cover page: title
      in uppercase, ``by'' (or its translation) and author name, and an
      address block containing agent and author contact information.
      The page counter is reset to~1 afterwards.
  \item[] \verb|\nicholl|\quad  Produces a simpler title page variant
      with just the title centred on the page and no address block.  The
      page counter is also reset to~1.
\end{description}

\subsection{Sluglines (scene headings)}

Sluglines mark the beginning of scenes, typically with ``INT.'' or
``EXT.'', the location, and time of day.  The class offers commands that
build upon the language-dependent abbreviations defined in the language
files.
\begin{description}
  \item[] \verb|\intslug[⟨qualifier⟩]{⟨location⟩}|\quad  Interior
    slugline.  The class prints the language-specific ``INT''
    abbreviation, followed by the location and any qualifier (such as
    ``DAY'', ``NIGHT'').  The optional \verb|⟨qualifier⟩| argument is
    placed at the end of the slugline, separated by a symbol, usually
    a dash.
  \item[] \verb|\extslug[⟨qualifier⟩]{⟨location⟩}|\quad  Exterior
    slugline.  As above, but using the ``EXT'' abbreviation.
  \item[] \verb|\intextslug[⟨qualifier⟩]{⟨location⟩}|\quad  Combined
    interior/exterior slugline (``INT./EXT.''), useful for scenes that
    move between inside and outside.
  \item[] \verb|\extintslug[⟨qualifier⟩]{⟨location⟩}|\quad  Combined
    exterior/interior slugline (``EXT./INT.'').
\end{description}

In \texttt{literary} mode, sluglines are printed without scene numbers.
In \texttt{technical} mode, each slugline increments an internal scene
counter and prints the scene number in the margins.

A typical usage might be:
\begin{verbatim}
    \intslug{OFFICE -- DAY}
    \extslug[NIGHT]{STREET}
\end{verbatim}

\subsection{Dialogue}

Dialogue is set using a dedicated environment which handles character
names, optional parenthetical remarks, and the width and indentation of
dialogue blocks.

\begin{description}
  \item[] \verb|\begin{dialogue}[⟨paren⟩]{⟨CHARACTER⟩} ... \end{dialogue}|\quad
    Starts a dialogue block for the given character name.  The character
    name is printed in uppercase at the standard dialogue position.  The
    optional \verb|⟨paren⟩| argument allows you to add a parenthetical
    remark (such as ``whispering''), which is typeset in parentheses
    above the dialogue text.

    Example:
\begin{verbatim}
    \begin{dialogue}[whispering]{WRITER}
        This is a whispered line.
    \end{dialogue}
\end{verbatim}

  \item[] \verb|\dialbreak[⟨paren⟩]{⟨CHARACTER⟩}|\quad  Ends the current
    \verb|dialogue| environment with a ``(MORE)'' marker, starts a new
    page, and opens a new dialogue block for the same character,
    labelled with a continuation marker (typically ``(CONT'D)'' or its
    translation).  The optional \verb|⟨paren⟩| argument works as for
    \verb|dialogue|.

    This is useful when a character's speech continues across a page
    break in the final layout.
\end{description}

Parenthetical remarks can also be inserted manually using:
\begin{description}
  \item[] \verb|\paren{⟨text⟩}|\quad  Typesets a parenthetical remark in
    a narrower column, typically centred under the character name and
    above the next line of dialogue.
\end{description}

\subsection{Transitions and labels}

The class provides a number of commands for common transitions and
on-screen text labels.  Their exact spelling and punctuation are
language-dependent and defined in the language files.
\begin{description}
  \item[] \verb|\fadein|\quad  Inserts the language-specific ``FADE
    IN:'' label followed by a vertical skip.  This is typically used at
    the beginning of the script.
  \item[] \verb|\fadeout|\quad  Inserts the language-specific ``FADE
    OUT:'' label (or its equivalent), usually aligned to the right in
    the standard transition position.
  \item[] \verb|\intercut|\quad  Inserts an ``INTERCUT WITH:'' label,
    usually flush right, to indicate intercutting between locations.
  \item[] \verb|\flashback|\quad  Inserts a ``FLASHBACK TO:'' label (or
    equivalent), again in the standard transition position.
\end{description}

For on-screen textual titles, use the \verb|titleover| environment:
\begin{description}
  \item[] \verb|\begin{titleover}[⟨qualifier⟩] ... \end{titleover}|\quad
    Starts a ``TITLE OVER:'' block.  The optional \verb|⟨qualifier⟩|
    argument allows additional specification (for example, a location or
    person).  The body of the environment is typeset as the on-screen
    text.

    Within this environment, \verb|\titbreak| can be used to mark a
    break to the next page while indicating that the title continues.
\end{description}

Finally, to mark the end of the script:
\begin{description}
  \item[] \verb|\theend|\quad  Typesets the language-specific ``THE
    END'' (or equivalent) centred on the page.
\end{description}

%}}}

\section{Language files}%{{{

All language-dependent strings used by the \texttt{screenwriter} class
are defined in separate language files named
\verb|screenwriter-lang-⟨language⟩.ldf|. This keeps the core class code
independent of any particular language and makes it easy to add new
languages or adjust existing ones.

\subsection*{Structure of a language file}

A typical language file has the following structure:

\begin{verbatim}
\ProvidesFile{screenwriter-lang-spanish.ldf}[2026/05/09 v0.1 Spanish]

\renewcommand*{\myaddresstext}{Agencia de contacto}
\renewcommand*{\mymoretext}{(M\'AS)}
\renewcommand*{\mycontdtext}{(CONT.)}
\renewcommand*{\myoffscreentext}{(F.P.)}
\renewcommand*{\myvoiceovertext}{(\emph{OFF})}
\renewcommand*{\myfadeintext}{ABRE DE NEGRO}
\renewcommand*{\myfadeouttext}{FUNDE A NEGRO}
\renewcommand*{\myadlibtext}{\emph{ad~lib.}}
\renewcommand*{\mypunctcharA}{:}
\renewcommand*{\mypunctcharB}{.}
\renewcommand*{\myplacesep}{.~}
\renewcommand*{\mysepintext}{./}
\renewcommand*{\mybytext}{por}
\renewcommand*{\myinttext}{INT}
\renewcommand*{\myexttext}{EXT}
\renewcommand*{\mytitleovertext}{R\'OTULO}
\renewcommand*{\myintercuttext}{CORTE A}
\renewcommand*{\myflashbacktext}{ANALEPSIS A}
\renewcommand*{\mypovtext}{P.D.V.}
\renewcommand*{\myreverttext}{REGRESO DE \pov}
\renewcommand*{\mythirtytext}{FIN}
\renewcommand*{\myslugspace}{ -- }
\end{verbatim}

The file starts with a \verb|\ProvidesFile| declaration, followed by a
series of \verb|\renewcommand*| statements that assign the language-
specific strings to the internal macros used by the class. No packages
are loaded and no options are processed in these files; they only
contain definitions.

The class expects the following macros to be defined (or redefined) by
each language file:
\begin{itemize}
  \item Cover and address strings:\\
    \verb|\myaddresstext|, \verb|\mybytext|.
  \item Dialogue markers:\\
    \verb|\mymoretext|, \verb|\mycontdtext|,
    \verb|\myoffscreentext|, \verb|\myvoiceovertext|,
    \verb|\myadlibtext|.
  \item Transitions and labels:\\
    \verb|\myfadeintext|, \verb|\myfadeouttext|,
    \verb|\mytitleovertext|, \verb|\myintercuttext|,
    \verb|\myflashbacktext|.
  \item Scene heading components:\\
    \verb|\myinttext|, \verb|\myexttext|,
    \verb|\myplacesep|, \verb|\mysepintext|,
    \verb|\myslugspace|.
  \item Other:\\
    \verb|\mypovtext|, \verb|\myreverttext|,
    \verb|\mythirtytext|.
\end{itemize}

\subsection*{Customizing individual strings}

In some cases, users may wish to keep a given language but customise the
wording of one or two internal strings (for example, preferring ``SMASH
CUT TO:'' instead of ``CUT TO:''). Since all language-dependent text is
stored in macros, this can be done directly in the document preamble.

For example, to change the ``FADE IN'' label while still using the
English language file:
\begin{verbatim}
\documentclass[english]{screenwriter}

% custom wording for the opening transition
\renewcommand*{\myfadeintext}{SMASH CUT IN}

\begin{document}
    \fadein
    ...
\end{document}
\end{verbatim}

Likewise, to adjust only the ``THE END'' marker in Spanish:

\begin{verbatim}
\documentclass[spanish]{screenwriter}

% use a different final marker
\renewcommand*{\mythirtytext}{FIN DE LA OBRA}

\begin{document}
    ...
    \theend
\end{document}
\end{verbatim}

When using \texttt{babel}, these customisations should be placed in the
preamble \emph{after} loading the class and any multilingual packages,
so that they override the values set by the language definition file.

\subsection*{Adding a new language}

To add support for a new language, you only need to:
\begin{enumerate}
  \item Create a new file
    \verb|screenwriter-lang-⟨language⟩.ldf| in the \texttt{lang/}
    subdirectory, based on one of the existing language files.
  \item Provide appropriate translations for all \verb|\my...| macros.
  \item Optionally, define a class option in \texttt{screenwriter.cls}
    so users can select the language explicitly, e.g.
\begin{verbatim}
\DeclareOption{⟨language⟩}{%
  \def\screenwriter@lang{⟨language⟩}%
  \screenwriter@lang@explicittrue
}
\end{verbatim}
\end{enumerate}

If no explicit class option is given, but the document uses
\texttt{babel} and the current \verb|\languagename| matches the file
name, \texttt{screenwriter} will automatically pick up the new language
file without further changes to the class.

%}}}

\section{License and credits}%{{{

% LPPL, agradecimiento a John Pate, etc.
This work may be distributed and/or modified under the conditions of the
\emph{\LaTeX{} Project Public License} (LPPL), either version~1.3c of
this license or (at your option) any later version.

The latest version of the LPPL is available at:
\begin{center}
  \url{https://www.latex-project.org/lppl.txt}
\end{center}

\noindent
Feedback, bug reports, and suggestions for new language files are very
welcome.

\subsection*{Credits}

The design of \textsf{screenwriter} is based on John Pate's original
\texttt{screenplay} class, which provided an early solution for
typesetting screenplays in \LaTeX{}.

The present class updates and extends that work for current \LaTeX{}
distributions, introduces modular language support, and refines several
aspects of the layout and implementation. Any remaining mistakes or
omissions are the responsibility of the current maintainer.

%}}}

\end{document}

%}}}
